\begin{enumerate}
    \item Les ondes ultrasonores sont longitudinales, car la direction de
          déformation du milieu de propagation est parallèle à la direction de
          propagation de l'onde.~\textbf{(1pt)}
    \item On trouve graphiquement~:
          $\tau=\qty{2}{\ms}\times\qty[per-mode=symbol]{2}{\ms\per\div}=
          \qty{4}{\ms}$.~\textbf{(1pt)}
    \item Le pétrole est plus dense que l'air, donc~: $v_{p} > v_{\text{air}}$\\
          En comparant les durées mises par les deux ondes sonores entre
          émetteur et le récepteur, on écrit~: $t_{p}<t_{\text{air}}$ \\
          Le retard temporel entre les deux ondes est~:
          $\tau=t_{\text{air}}-t_{p}$~; avec
          $t_{\text{air}}=\frac{L}{v_{\text{air}}}$ et $t_{p}=\frac{L}{v_{p}}$.
          
          Finalement~:~\textbf{(1pt)}
          \[
              \tau=L \cdot\left(\frac{1}{v_{\text{air}}}-\frac{1}{v_p}\right)
          \]
    \item $\tau=L \cdot\left(\frac{1}{v_{\text{air}}}-\frac{1}{v_p}\right)
          \quad\Rightarrow\
          \frac{\tau}{L}=\frac{1}{v_{\text{air}}}-\frac{1}{v_p}
          \quad\Rightarrow\
          \frac{1}{v_p}=\frac{1}{v_{\text{air}}}-\frac{\tau}{L}=
          \frac{L-\tau v_{\text{air}}}{Lv_{\text{air}}}
          \quad\Rightarrow\
          v_{p}=\frac{Lv_{\text{air}}}{L-\tau v_{\text{air}}}$~\textbf{(0,5pt)}

          A.N. $v_{p}=\frac{1,84\times 340}{1,84-\num{4e-3}\times 340}$
          \quad $\Rightarrow\ v_{p}=\qty{1303}{\v}$
\end{enumerate}

